\section{$(\mu,\epsilon)$-Robustness}
\label{sec:new-robustness}

We define classifier robustness in its usual form, survey related work containing other approaches to defining global robustness and highlight differences to our proposed formulation.

\paragraph{Local Robustness.} Local classifier robustness is an intuitive widely accepted in the literature property, which we formulate over semimetric spaces. We use the notation $[n] := \{1, 2, \dots,n\} \subset \mathbb{N}$.
% lax begin Lax544225.Classifier
\begin{definition}[Classifier]
\label{def:classifier}
Let $S$ be a set of \enquote{inputs}  and  $[n]$ a set of \enquote{classes}. 
A \emph{classifier} is a total, surjective function:
$c: S \rightarrow [n]$. 
\end{definition}
% lax end

Surjectivity ensures that each class is represented by at least one input, reflecting the common requirement in machine learning that all labels occur in the domain and helping to avoid systematic bias toward missing classes.
% lax begin Lax544225.Semimetric
\begin{definition}[Semimetric Space]
\label{def:distance}
    Let $(S, d)$ be a pair of a nonempty set and a function $d: S \times S \rightarrow \mathbb{R}_{\geq 0}$ such that for all $x,y \in S$, $d(x,y) = 0 \iff x = y$ (identity of indiscernibles) and $d(x,y) = d(y,x)$ (symmetry). Then, we call $(S, d)$ a \emph{semimetric space} and $d$ a \emph{distance} function on $S$.
\end{definition}
% lax end

Using a distance, we define local $\epsilon$-robustness at an input $x^*$ requiring that sufficiently close (similar) inputs receive the same label. If it holds for all possible inputs, the classifier is globally $\epsilon$-robust.
% lax begin Lax544225.ClassifierRobustness
\begin{definition}[$\epsilon$-Robustness]
\label{def:local-robustness}
Let $(S,d)$ be a semimetric space. A classifier $c: S \rightarrow [n]$ is \emph{locally robust} at a point $x^* \in S$ for $\epsilon > 0$ if:
$$\forall x \in S [d(x, x^*) \leq \epsilon \rightarrow c(x) = c(x^*)].$$    
\end{definition}
% lax end

If a classifier $c$ is given, we call $x \in S$ a \textit{(non-)robust point} if $c$ is (not) locally robust on $x$. If a classifier is robust at all points in $S$, it is \emph{globally} $\epsilon$-robust.

\paragraph{Research Gap.} Consider the question of defining global robustness. An obvious way is to enforce local robustness for every point. But as ref. \cite{li2024certifying} points out, \textit{"local robustness verification cannot be directly generalized to the whole decision region of a class"} because the distance can \textit{"cross decision boundaries [...], making a verification task unrealistic"}. It is easy to see that points close to the true decision boundary are not robust which makes robustness on every point infeasible. 

Most works explore global robustness under additional assumptions. Different authors restrict robustness to specific datasets or underlying data distributions~\cite{casadio2022neural,tsipras2018robustness,zhang2019theoretically,yang2020closer,mangal2023certifying}. Some authors study the trade-off between robustness and accuracy and demonstrate that, for a particular classification problem, improvements in accuracy necessarily come at the cost of robustness~\cite{tsipras2018robustness,zhang2019theoretically}.
In contrast, others show that robustness and accuracy can positively correlate when class distributions are separable, a property observed in several standard image datasets~\cite{yang2020closer,mangal2023certifying}. 

\begin{figure}[t]
\centering
\definecolor{muBlue}{HTML}{1F4E79}
\definecolor{epsOrange}{HTML}{D9730D}
\begin{tikzpicture}[scale=0.58]
  % lightly shaded l-infinity epsilon-balls of the minority inputs:
  % together they cover the whole grid
  \foreach \x in {2,7}{
    \foreach \y in {2,7}{
      \fill[orange!10] (\x-2.5,\y-2.5) rectangle (\x+2.5,\y+2.5);
    }
  }

  % faint grid to read the lattice as a grid of inputs;
  % the thick ball boundaries run between rows and columns of inputs
  \foreach \i in {0,...,9}{\draw[gray!18] (\i,-0.5)--(\i,9.5);}
  \foreach \j in {0,...,9}{\draw[gray!18] (-0.5,\j)--(9.5,\j);}

  % l-infinity epsilon-balls of the minority inputs (half-width eps=2.5)
  \foreach \x in {2,7}{
    \foreach \y in {2,7}{
      \draw[epsOrange, very thick] (\x-2.5,\y-2.5) rectangle (\x+2.5,\y+2.5);
    }
  }

  % eps arrow (half-width of an l-inf ball), drawn beneath the inputs
  \draw[<->,thick,gray!70] (2,2) -- (4.5,2);
  \node[anchor=south,inner sep=1.5pt] at (3.3,2.05) {\footnotesize $\epsilon$};

  % majority-label (class A) inputs on the whole grid
  \foreach \i in {0,...,9}{
    \foreach \j in {0,...,9}{
      \fill[muBlue] (\i,\j) circle (1.7pt);
    }
  }

  % minority-label (class B) inputs on a sparse sublattice
  \foreach \x in {2,7}{
    \foreach \y in {2,7}{
      \fill[epsOrange] (\x-0.17,\y) -- (\x,\y+0.17) -- (\x+0.17,\y) -- (\x,\y-0.17) -- cycle;
    }
  }

  % compact legend to the right of the grid
  \begin{scope}[shift={(10.7,5.7)}]
    \fill[muBlue] (0,0) circle (1.7pt);
    \node[anchor=west] at (0.35,0) {\footnotesize majority label $A$};
    \fill[epsOrange] (-0.17,-1) -- (0,-0.83) -- (0.17,-1) -- (0,-1.17) -- cycle;
    \node[anchor=west] at (0.35,-1) {\footnotesize minority label $B$};
    \draw[epsOrange, very thick, fill=orange!10] (-0.05,-2.5) rectangle (0.5,-1.95);
    \node[anchor=west] at (0.7,-2.225) {\footnotesize $\epsilon$-ball of $B$};
  \end{scope}
\end{tikzpicture}
\caption{Probabilistic robustness does not constrain the \emph{structure} of
non-robust points. In the image, inputs of class $B$ are rare, so a randomly drawn
$\epsilon$-close pair is almost always same-label and density-based
robustness bound is met. Yet the $B$ inputs sit on a sublattice whose
$\epsilon$-balls (which become squares with $\ell_\infty$ norm) cover the entire input domain: every input of
the majority class $A$ lies within $\epsilon$ of some $B$, so no input is locally robust.}
\label{fig:grid}
\end{figure}

When it comes to measuring global robustness, authors introduce additional concepts to avoid working with robustness at every point. A common approach is to enforce robustness only on points where a classifier has high confidence~\cite{athavale2024verifying,leino2021globally}. This solution creates another problem of ensuring that classifier prediction does reach high confidence whenever required~\cite{el2010foundations,chow1970optimum}. Another popular approach is to enforce that a probability of finding a pair $x, x^*\in S$ with $d(x,x^*) \leq \epsilon$ and $c(x) \neq c(x^*)$ is less than some threshold value~\cite{levy2023groma,li2024certifying,mangal2019robustness}. Probabilistic approach only bounds the density of non-robust points and does not tell anything about their structure. For example, a grid of singular points such as one in Figure~\ref{fig:grid} with a different classifier label from its neighbors can result in a satisfaction of probabilistic robustness for a classifier that has no robust points (as they have a non-robust point from the grid as a neighbor).

\paragraph{This Work: $(\mu,\epsilon)$-Robustness.} In this work, we explore a new version of global classifier robustness. We consider a classifier to be $(\mu,\epsilon)$-robust if every non-robust point has a robust point located at most $\mu$-far. Note that $(0,\epsilon)$-robustness is equivalent to $\epsilon$-robustness.
% lax begin Lax544225.ClassifierRobustness
\begin{definition}[$(\mu,\epsilon)$-Robustness]
\label{def:mueps}
Let $(S,d)$ be a semimetric space and\linebreak $c: S \rightarrow [n]$ a classifier. $c$ is $(\mu, \epsilon)$-robust for $\mu \in \mathbb{R}_{\geq0}$ and $\epsilon \in \mathbb{R}_{>0}$ if for all $x \in S$, there exists $x' \in S$ with $d(x,x') \leq \mu$ such that $c(x) = c(x')$ and $c$ is locally $\epsilon$-robust at $x'$.
\end{definition}
% lax end

This property is a novel domain-independent version of global classifier robustness. Uniquely, this property enforces additional structure on non-robust points, enforcing coverage of the input domain by robust points. In the following sections, we will explore properties of $(\mu,\epsilon)$-robust classifiers in a closer detail.